\section{Các quy tắc tính đạo hàm}

\subsection{Đề 1}

\Opensolutionfile{ans}[ans/ans-K11_C7_B2]
\caulc

\begin{ex}%[1D7H2-1]%Câu 1
	Cho hàm số $ y=\dfrac{4}{x-1}$. Khi đó $y'\left(-1\right)$bằng
	\choice
	{\True $-1$}
	{$-2$}
	{$2$}
	{$1$}
	\loigiai{
				Ta có $y'=-\dfrac{4}{\left(x-1\right)^2}$ $\Rightarrow{y}'\left(-1\right)=-1$.}
\end{ex}

\begin{ex}%Câu 2%[1D7H2-1]
	Tính đạo hàm của hàm số $y=\sqrt{x}+x$ tại điểm $x_0=4$ là:
	\choice
	{$y'(4)=\dfrac{9}{2}$}
	{$y'(4)=6$}
	{$y'(4)=\dfrac{3}{2}$}
	{\True $y'(4)=\dfrac{5}{4}$}
	\loigiai{	
		Ta có $y'=\dfrac{1}{2\sqrt{x}}+1$ $\Rightarrow{y}'(4)=\dfrac{1}{2\sqrt{4}}+1=\dfrac{5}{4}$.}
\end{ex}

\begin{ex}%Câu 3%[1D7N3-1]
	Cho hàm số $ f(x)=x^3+2x$, giá trị của $f''(1)$ bằng
	\choice
	{\True $ 6$}
	{$ 8$}
	{$ 3$}
	{$ 2$}
	\loigiai{
			$f'(x)=3x^2+2$, $f''(x)=6x \Rightarrow  f''(1)=6$.}
\end{ex}

\begin{ex}%Câu 4%[1D7N3-1]
	Cho hàm số $ y=x^5-3x^4+x+1$ với $ x\in\mathbb{R}$. Đạo hàm $y''$ của hàm số là
	\choice
	{$y''=5x^3-12x^2+1$}
	{$y''=5x^4-12x^3$}
	{$y''=20x^2-36x^3$}
	{\True $y''=20x^3-36x^2$}
	\loigiai{
			Ta có $ y=x^5-3x^4+x+1$$\Rightarrow{y}'=5x^4-12x^3+1\Rightarrow y''=20x^3-36x^2$.}
\end{ex}

\begin{ex}%Câu 5%[1D7H2-1]
	Tính đạo hàm của hàm số $ f(x)=\dfrac{2x}{x-1}$ tại điểm $ x=-1$.
	\choice
	{$ f'\left(-1\right)=1$}
	{\True $ f'\left(-1\right)=-\dfrac{1}{2}$}
	{$ f'\left(-1\right)=-2$}
	{$ f'\left(-1\right)=0$}
	\loigiai{
			Ta có $ f'(x)=\dfrac{-2}{\left(x-1\right)^2}$. Vậy $ f'\left(-1\right)=-\dfrac{1}{2}$.}
\end{ex}

\begin{ex}%Câu 6%[1D7H2-1]
	Cho hàm số $y=x^2-x+2$. Tính $y'(1)$ .
	\choice
	{$y'(1)=-1$}
	{\True $y'(1)=1$}
	{$y'(1)=2$}
	{$ y'(1)=0$}
	\loigiai{
			$y'=2x-1$ $\Rightarrow y'(1)=2\cdot 1-1=1$.}
\end{ex}

\begin{ex}%Câu 7%[1D7H3-1]
	Cho $f(x)=x^3$. Tính $f''(1)$.
	\choice
	{$f''(1)=3$}
	{$f''(1)=2$}
	{\True $f''(1)=6$}
	{$f''(1)=1$}
	\loigiai{
		$f(x)=x^3\Rightarrow{f}'(x)=3x^2\Rightarrow f''(x)=3.2x=6x \Rightarrow f''(1)=6\cdot 1=6$}
\end{ex}

\begin{ex}%Câu 8%[1D7H3-2]
	Cho $ y=\sqrt{2x-x^2}$, tính giá trị biểu thức $ A=y^3\cdot y''$.
	\choice
	{$ 1$}
	{$ 0$}
	{\True $-1$}
	{$4$}
	\loigiai{	
		Ta có $ y'=\dfrac{1-x}{\sqrt{2x-x^2}}$, $y''=\dfrac{-1}{\left(\sqrt{2x-x^2}\right)^3}$.\\
		Do đó $ A=y^3.y''=-1$.}
\end{ex}

\begin{ex}%Câu 9%[1D7H3-1]
	Cho hàm số $ f(x)=x^3+2x$, giá trị của $f''(1)$ bằng
	\choice
	{\True $ 6$}
	{$ 8$}
	{$ 3$}
	{$ 2$}
	\loigiai{
			$f'(x)=3x^2+2$, $f''(x)=6x \Rightarrow f''(1)=6$.}
\end{ex}

\begin{ex}%Câu 10%[1D7H2-1]
	Cho hàm số $ f(x)=\sqrt{x^2-2 x+3}$, tính $f'(2)$.
	\choice
	{$\dfrac{1}{3}$}
	{$\sqrt{3}$}
	{\True $\dfrac{\sqrt{3}}{3}$}
	{$ 2\sqrt{3}$}
	\loigiai{
	Ta có $ f(x)=\sqrt{x^2-2x+3}\Rightarrow{f}'(x)=\dfrac{x-1}{\sqrt{x^2-2x+3}}\Rightarrow{f}'(2)=\dfrac{\sqrt{3}}{3}$.}
\end{ex}

\begin{ex}%Câu 11%[1D7H3-1]
	Cho hàm số $ y=x^5-3x^4+x+1$ với $ x\in\mathbb{R}$. Đạo hàm $y''$ của hàm số là
	\choice
	{$y''=5x^3-12x^2+1$}
	{$y''=5x^4-12x^3$}
	{$y''=20x^2-36x^3$}
	{\True $y''=20x^3-36x^2$}
	\loigiai{
		Ta có $ y=x^5-3x^4+x+1$$\Rightarrow{y}'=5x^4-12x^3+1\Rightarrow y''=20x^3-36x^2$.}
\end{ex}

\begin{ex}%Câu 12%[1D7H3-1]
	Đạo hàm cấp hai của hàm số $y=\cos^2x$ là
	\choice
	{\True $y''=-2\cos 2x$}
	{$y''=-2\sin 2x$}
	{$y''=2\cos 2x$}
	{$y''=2\sin 2x$}
	\loigiai{	
		$y'=2\cos x\cdot \left(-\sin x\right) =-\sin 2x \Rightarrow y''=-2\cos 2x$.}
\end{ex}

\Closesolutionfile{ans}
\indapan{6}{ans/ans-K11_C7_B2}
\cauds
\Opensolutionfile{ans}[ans/ans-K11_C7_B2-DS]

\begin{ex}%Câu 13%[1D7H2-3]
	Cho hàm số $y=x^4-4 x^2+3\sqrt{x}+2-\dfrac{1}{x}$. Khi đó
	\choiceTF
	{\True $ y'(1)=-\dfrac{3}{2}$}
	{Đồ thị của hàm số $ y'$ đi qua điểm $ A\left(1;\dfrac{3}{2}\right)$}
	{\True $y'(4)=\dfrac{3597}{16}$}
	{\True Điểm $ M$ thuộc đồ thị $(C)$của hàm số $y=x^4-4 x^2+3\sqrt{x}+2-\dfrac{1}{x}$ có hoành độ $x_0=1$. Khi đó, phương trình tiếp tuyến của $(C)$ tại $M$ vuông góc với đường thẳng $ y=\dfrac{2}{3}x$}
\loigiai{
$y^{\prime}=4\cdot x^3-4\cdot 2\cdot x+3\cdot\dfrac{1}{2\sqrt{x}}+0-\left(-\dfrac{1}{x^2}\right)=4 x^3-8 x+\dfrac{3}{2\sqrt{x}}+\dfrac{1}{x^2}$.
	\begin{itemchoice}
	\itemch Đúng. Vì $ y'(1)=-\dfrac{3}{2}$.
	\itemch Sai. Vì $ y'(1)=-\dfrac{3}{2}$, nên đồ thị của hàm số $ y'$ không đi qua điểm $ A\left(1;\dfrac{3}{2}\right)$.
	\itemch Đúng. Vì $y'(4)=\dfrac{3597}{16}$.
	\itemch Đúng. Hệ số góc tiếp tuyến tại $M$ là $k= y'(1)=-\dfrac{3}{2}$ và có $-\dfrac{3}{2}\cdot \dfrac{2}{3} = -1$ nên tiếp tuyến của $(C)$ tại $M$ vuông góc với đường thẳng $ y=\dfrac{2}{3}x$.
\end{itemchoice}
}
\end{ex}

\begin{ex}%Câu 14%[1D7V2-1]
Cho hàm số $y=\dfrac{x^2+1}{x-3}$. Khi đó
\choiceTF
{$ y'(1)=\dfrac{3}{2}$}
{Tổng các nghiệm của phương trình $ y'=0$ bằng $-6$}
{\True Đồ thị của hàm số $ y'$ đi qua điểm $ A\left(1;-\dfrac{3}{2}\right)$}
{$ y'(1)<y'(2)$}
\loigiai{
$y^{\prime}=\dfrac{\left(x^2+1\right)^{\prime}(x-3)-(x-3)^{\prime}\left(x^2+1\right)}{(x-3)^2}=\dfrac{2 x(x-3)-1\left(x^2+1\right)}{(x-3)^2}=\dfrac{x^2-6 x-1}{(x-3)^2}$.
	\begin{itemchoice}
	\itemch Sai. Vì $ y'(1)=-\dfrac{3}{2}$.
	\itemch Sai. Tổng các nghiệm của phương trình $ y'=0$ bằng $6$.
	\itemch Đúng. Vì $ y'(1)=-\dfrac{3}{2}$.
	\itemch Sai. $y'(1) = -1 > y'(2) = -9$.
\end{itemchoice}
}
\end{ex}

\begin{ex}%Câu 15%[1D7H2-1]
Tính được đạo hàm của các hàm số sau. Khi đó
\choiceTF
{\True $y=\log_2(9 x-5)$ có $y^{\prime}=\dfrac{9}{(9x-5)\ln 2}$}
{\True $y=2 e^{3 x+1}$ có $y^{\prime}=6e^{3x+1}$}
{\True $y=3^{x^3-1}$ có $y^{\prime}=3\cdot\ln 3\cdot{x^2}\cdot{3^{x^3-1}}$}
{$y=\ln\sqrt{x}$ có $y^{\prime}=-\dfrac{1}{2x}$}
\loigiai{
	\begin{itemchoice}
	\itemch Đúng. Vì $y^{\prime}=\dfrac{(9 x-5)^{\prime}}{(9 x-5)\ln 2}=\dfrac{9}{(9 x-5)\ln 2}$.
	\itemch  Đúng. Vì $y^{\prime}=2(3 x+1)^{\prime}\cdot e^{3 x+1}=6 e^{3 x+1}$.
	\itemch Đúng. Vì $y^{\prime}=\left(x^3-1\right)^{\prime}\cdot 3^{x^3-1}\cdot\ln 3=3\cdot\ln 3\cdot x^2\cdot 3^{x^3-1}$.
	\itemch Sai. Vì $y^{\prime}=\dfrac{(\sqrt{x})^{\prime}}{\sqrt{x}}=\dfrac{1}{2\sqrt{x}\sqrt{x}}=\dfrac{1}{2 x}$.
\end{itemchoice}}
\end{ex}

\begin{ex}%Câu 16%[1D7V3-3]
Chuyển động của một vật có phương trình $s(t)=4\cos\left(2\pi t-\dfrac{\pi}{12}\right)$ m, với $t$ là thời gian tính bằng giây. Khi đó
\choiceTF
{\True $s^{\prime}(t)=-8\pi\sin\left(2\pi t-\dfrac{\pi}{12}\right)$}
{$s^{\prime\prime}(t)=16\pi^2\cos\left(2\pi t-\dfrac{\pi}{12}\right)$}
{\True Vận tốc của vật tại thời điểm khi $t=5$ s là 
$s'(5)= \text{ m/s}$}
{Gia tốc của vật tại thời điểm khi $t=5$ s là $s''(5)\approx 152,533$ m/$s^2$}
\loigiai{
		\begin{itemchoice}
		\itemch  Đúng. Vì ta có $s^{\prime}(t)=-8\pi\sin\left(2\pi t-\dfrac{\pi}{12}\right)$ và $s^{\prime\prime}(t)=-16\pi^2\cos\left(2\pi t-\dfrac{\pi}{12}\right)$.
		\itemch  Sai. Vì ta có $s^{\prime}(t)=-8\pi\sin\left(2\pi t-\dfrac{\pi}{12}\right)$ và $s^{\prime\prime}(t)=-16\pi^2\cos\left(2\pi t-\dfrac{\pi}{12}\right)$.
		\itemch  Đúng. Vì vận tốc tức thời của vật tại thời điểm $t=5$ s là \\
$s^{\prime}(5)=-8\pi\sin\left(10\pi-\dfrac{\pi}{12}\right)\approx 6,505$ (m/s).
	\itemch   Đúng. Vì gia tốc tức thời của vật tại thời điểm $t=5$ s là\\
$s^{\prime\prime}(5)=-16\pi^2\cos\left(10\pi-\dfrac{\pi}{12}\right)\approx-152{,}533$ (m/s$^2$).
\end{itemchoice}
}
\end{ex}


\Closesolutionfile{ans}
\indapan{3}{ans/ans-K11_C7_B2-DS}

\Opensolutionfile{ans}[ans/ans-K11_C7_B2-KQ]
\caukq

\begin{ex}%Câu 17%[1D7V2-8]
	Dân số (tính theo nghìn người) của một thành phố được cho bởi công thức $f(t)=\dfrac{26 t+10}{t+5}$, trong đó $t$ (được tính bằng năm) là khoảng thời gian kể từ năm $2015$. Tìm tốc độ tăng dân số trong năm $2025$ của thành phố đó (làm tròn đến phần trăm).
\shortans{$ 0{,}53$}
\loigiai{
	Đạo hàm của hàm số $f$ biểu thị tốc độ tăng dân số của thành phố đó (tính bằng nghìn người/ năm), ta có $f^{\prime}(t)=\dfrac{120}{(t+5)^2}$.\\
	Từ năm 2015 đến năm 2025 nghĩa là $t=10$.\\
	Vậy tốc độ tăng dân số tại thời điểm $t=10$ là\\
	$f^{\prime}(10)=\dfrac{120}{(10+5)^2}=\dfrac{8}{15}\approx 0{,}53$ (nghìn người/năm)
}
\end{ex}

\begin{ex}%Câu 18%[1D7V2-8]
Cân nặng trung bình của một em bé trong độ tuổi từ 0 đến 36 tháng có thể được tính gần đúng bởi hàm số $w(t)=0,00076 t^3-0,06 t^2+1,8 t+8,2$, trong đó $t$ được tính bằng tháng và $w$ được tính bằng pound. Tính tốc độ thay đổi cân nặng của em bé đó tại thời điểm 15 tháng tuổi.\\
\shortans{$0{,}51$}
\loigiai{
Ta có $w^{\prime}(t)=\dfrac{57}{25000}t^2-\dfrac{3}{25}t+1{,}8$.\\
Tốc độ thay đổi cân nặng của em bé đó tại thời điểm $15$ tháng tuổi là\\
$w^{\prime}(15)=\dfrac{57}{25000}\cdot 15^2-\dfrac{3}{25}\cdot 15+1{,}8=0{,}51$ (pound/tháng)
}
\end{ex}

\begin{ex}%Câu 19%[1D7H2-1]
Đạo hàm của hàm số $y=\dfrac{\sin ^2 x-\cos ^2 x}{\sin x\cdot\cos x}$ tại điểm $x=\dfrac{\pi}{6}$ (làm tròn đến phần trăm).
\shortans{$5{,}33$}
\loigiai{
$y=\dfrac{\sin ^2 x-\cos ^2 x}{\sin x\cdot\cos x}=\dfrac{-\cos 2 x}{\dfrac{1}{2}\sin 2 x}=-2\cot 2 x$\\
$\Rightarrow y^{\prime}=-2\dfrac{-2}{\sin ^2 2 x}=\dfrac{4}{\sin ^2 2 x}\Rightarrow y^{\prime}\left(\dfrac{\pi}{6}\right)=\dfrac{16}{3}\approx 5{,}33$.}
\end{ex}

\begin{ex}%Câu 20%[1D7H2-1]
Chứng minh hàm số $f(x)=\sin ^6 x+\cos ^6 x+3\sin ^2 x\cdot\cos ^2 x-2 x$ có đạo hàm không phụ thuộc $x$.\\
\shortans{$-2$}
\loigiai{
$\begin{aligned}
f^{\prime}(x)& =\left(\sin ^6 x+\cos ^6 x+3\sin ^2 x\cdot\cos ^2 x-2 x\right)^{\prime}\\
&=6\sin ^5 x\cdot\cos x-6\cos ^5 x\sin x+6\sin x\cdot\cos ^3 x-6\sin ^3 x\cdot\cos x-2\\
&=6\sin x\cos x\left(\sin ^4 x-\cos ^4 x+\cos ^2 x-\sin ^2 x\right)-2\\
&=6\sin x\cos x\left(\sin ^2 x\left(\sin ^2 x-1\right)+\cos ^2 x\left(1-\cos ^2 x\right)\right)-2\\
&=6\sin x\cos x\left(-\cos ^2 x\sin ^2 x+\cos ^2 x\sin ^2 x\right)-2=-2 .
\end{aligned}$}
\end{ex}

\begin{ex}%Câu 21%[1D7V1-5]
Một vật chuyển động theo quy luật $s=\dfrac{1}{3}t^3-t^2+9 t$, với $t$ (giây) là khoảng thời gian tính từ lúc vật bắt đầu chuyển động và $s$ (mét) là quãng đường vật đi được trong thời gian đó. Hỏi trong khoảng thời gian $10$ giây, kể từ lúc bắt đầu chuyển động, vận tốc nhỏ nhất của vật đạt được bằng bao nhiêu?\\
\shortans{ $8$}
\loigiai{
Vận tốc $v$ của vật được tính theo công thức $v(t)=s^{\prime}(t)=t^2-2 t+9$.\\
Ta có $t^2-2 t+9=(t-1)^2+8\geq 8\Rightarrow v\geq 8$.\\
Vậy vận tốc nhỏ nhất của vật là $8$ m/s đạt được tại thời điểm $t=1$ giây.}
\end{ex}

\begin{ex}%Câu 22%[1D7V1-5]
Cho chuyển động thẳng xác định bởi phương trình $S=-t^3+3 t^2+9 t$, trong đó $t$ tính bằng giây và $S$ tính bằng mét. Tính vận tốc của chuyển động tại thời điểm gia tốc triệt tiêu.\\
\shortans{ $ 12$}
\loigiai{
Ta có: $v=S^{\prime}=-3 t^2+6 t+9$ và $a=v^{\prime}=-6 t+6$.\\
Gia tốc triệt tiêu khi $a=0\Rightarrow t=1$.\\
Khi đó vận tốc của chuyển động là $v(1)=12$ m/s.
}
\end{ex}



\Closesolutionfile{ans}
\indapan{6}{ans/ans-K11_C7_B2-KQ}

\subsection{Đề 2}

\Opensolutionfile{ans}[ans/ans-K11_C7_B2]
\caulc

\begin{ex}%Câu 23%[1D7H2-1]
	Tính đạo hàm của hàm số $ f(x)=\dfrac{2x+7}{x+4}$ tại $ x=2$ ta được
	\choice
	{\True $f'(2)=\dfrac{1}{36}$}
	{$f'(2)=\dfrac{11}{6}$}
	{$f'(2)=\dfrac{3}{2}$}
	{$f'(2)=\dfrac{5}{12}$}
	\loigiai{	
		Ta có $f'(x)=\dfrac{1}{\left(x+4\right)^2}\Rightarrow f'(2)=\dfrac{1}{36}$.}
\end{ex}

\begin{ex}%Câu 24%[1D7H2-1]
	Cho hàm số $ f(x)$ xác định trên $\mathbb{R}$ bởi $ f(x)=2x^2+1$. Giá trị $f'\left(-1\right)$ bằng
	\choice
	{$2$}
	{$6$}
	{\True $-4$}
	{$3$}
	\loigiai{
			Ta có  $ f'(x)=4x$ $\Rightarrow{f}'\left(-1\right)=-4$.}
\end{ex}

\begin{ex}%Câu 25%[1D7H3-1]
	Đạo hàm cấp hai của hàm số $ y=x^6-4x^3+2x+2022$ với $ x\in\mathbb{R}$ là
	\choice
	{$y''=30x^4-24x+2$}
	{$y''=30x^4-24x$}
	{\True $y''=6x^5-12x^2+2$}
	{$y''=6x^5-12x^2$}
	\loigiai{	
		Ta có $y'=6x^5-12x^2+2$. 	Suy ra $y''=30x^4-24x$.}
			\end{ex}
\begin{ex}%[1D7H3-1]
		 Đạo hàm cấp hai của hàm số $ y=\cos^2x$ là
		\choice
		{\True $y''=-2\cos 2x$}
		{$y''=-2\sin 2x$}
		{$y''=2\cos 2x$}
		{$y''=2\sin 2x$}
		\loigiai{		
			$ y'=2\cos x\cdot \left(-\sin x\right) =-\sin 2x \Rightarrow y''=-2\cos 2x$.}
	\end{ex}
	
	\begin{ex}%Câu 26%[1D7H2-1]
		Tính đạo hàm của hàm số $y=x\left(x+1\right)\left(x+2\right)\left(x+3\right)$ tại điểm $x_0=0$ là
		\choice
		{$y'(0)=5$}
		{\True $y'(0)=6$}
		{$y'(0)=0$}
		{$y'(0)=-6$}
		\loigiai{
			Ta có $y=x\left(x+1\right)\left(x+2\right)\left(x+3\right)=\left(x^2+x\right)\left(x^2+5x+6\right)$\\
			$\Rightarrow{y}'=\left(2x+1\right)\left(x^2+5x+6\right)+\left(x^2+x\right)\left(2x+5\right)$ 
			$\Rightarrow{y}'(0)=6.$}
	\end{ex}
	

	
	\begin{ex}%Câu 28%[1D7H2-1]
		Cho hàm số $y=\dfrac{x+2}{x-1}$ . Tính $y'(3)$ 
		\choice
		{$\dfrac{5}{2}$}
		{\True $-\dfrac{3}{4}$}
		{$-\dfrac{3}{2}$}
		{$\dfrac{3}{4}$}
		\loigiai{
			Ta có $y=\dfrac{x+2}{x-1}\Rightarrow y'=\dfrac{-3}{\left(x-1\right)^2}$\\
			$y'(3)=\dfrac{-3}{\left(3-1\right)^2}=-\dfrac{3}{4}$.}
	\end{ex}
	
	\begin{ex}%Câu 29%[1D7H3-2]
		Cho hàm số $ y=x\cos x$. Tìm hệ thức đúng trong các hệ thức sau
		\choice
		{$y''+y=\sin x+2x\cos x$}
		{$y''+y=2\sin x$}
		{$y''+y=-\sin x+x\cos x$}
		{\True $y''+y=-2\sin x$}
		\loigiai{
			Ta có $y'=\text{cos}x-x\sin x\Rightarrow y''=-2\sin x-x\cos x$.\\
			Khi đó $y''+y=-2\sin x-x\cos x+x\cos x=-2\sin x$.}
	\end{ex}
	
	\begin{ex}%Câu 30%[1D7H3-1]
		Đạo hàm cấp hai của hàm số $ y=\cos^2x$ là
		\choice
		{\True $y''=-2\cos 2x$}
		{$y''=-2\sin 2x$}
		{$y''=2\cos 2x$}
		{$y''=2\sin 2x$}
		\loigiai{
				$ y'=2\cos x.\left(-\sin x\right)=-\sin 2x\Rightarrow y''=-2\cos 2x$.}
	\end{ex}
	
	\begin{ex}%Câu 31%[1D7H2-1]
		Cho hàm số $f(x)=\sqrt{x-1}$ . Đạo hàm của hàm số tại $ x=1$ là
		\choice
		{$\dfrac{1}{2}$}
		{$1$}
		{$0$}
		{\True Không tồn tại}
		\loigiai{
					Ta có $f'(x)=\dfrac{1}{2\sqrt{x-1}}$ do đó không tồn tại đạo hàm tại $x=1$.}
	\end{ex}
	
	\begin{ex}%Câu 32%[1D7H2-1]
		Cho hàm số $y=\dfrac{x^2+x}{x-2}$ . Đạo hàm của hàm số tại $x=1$ là
		\choice
		{$y'(1)=-4$}
		{\True $y'(1)=-5$}
		{$y'(1)=-3$}
		{$y'(1)=-2$}
		\loigiai{
				Ta có $y=x+3+\dfrac{6}{x-2} \Rightarrow y' = 1-\dfrac{6}{(x-2)^2}$. Vậy $y'(1)=-5.$}
	\end{ex}
	
	\begin{ex}%Câu 33%[1D7H2-1]
		Tính đạo hàm cấp hai của hàm số $ y=-3\cos x$ tại điểm $x_0=\dfrac{\pi}{2}$.
		\choice
		{$y''\left(\dfrac{\pi}{2}\right)=-3$}
		{$y''\left(\dfrac{\pi}{2}\right)=5$}
		{\True $y''\left(\dfrac{\pi}{2}\right)=0$}
		{$y''\left(\dfrac{\pi}{2}\right)=3$}
		\loigiai{		
			$ y=-3\cos x \Rightarrow{y}'=3\sin x \Rightarrow y''=3\cos x \Rightarrow y''\left(\dfrac{\pi}{2}\right)=0$.}
	\end{ex}
	
	\begin{ex}%Câu 34%[1D7H3-1]
		Đạo hàm cấp hai của hàm số $ y=\dfrac{3x+1}{x+2}$ là
		\choice
		{$y''=\dfrac{10}{\left(x+2\right)^2}$}
		{$y''=-\dfrac{5}{\left(x+2\right)^4}$}
		{$y''=-\dfrac{5}{\left(x+2\right)^3}$}
		{\True $y''=-\dfrac{10}{\left(x+2\right)^3}$}
		\loigiai{		
			Ta có $ y=3-\dfrac{5}{x+2}\Rightarrow{y}'=\dfrac{5}{\left(x+2\right)^2};{y}''=-\dfrac{10}{\left(x+2\right)^3}$\\
			}
	\end{ex}



\Closesolutionfile{ans}
\indapan{6}{ans/ans-K11_C7_B2}
\cauds
\Opensolutionfile{ans}[ans/ans-K11_C7_B2]

\begin{ex}%Câu 35%[1D7V2-1]
	Cho hàm số $y=-4 x^3+\dfrac{x^2}{2}-2 x+3$, biết $ y'=a{x^2}+bx+c$. Khi đó
	\choiceTF
	{$ a+b+c=-10$}
	{Phương trình $ y'=0$ có hai nghiệm phân biệt}
	{\True Đồ thị hàm số $ y'$ cắt trục tung tại điểm $\left(0;-2\right)$}
	{Đồ thị hàm số $ y'$ cắt đường thẳng $ y=3$ tại hai điểm phân biệt}
\loigiai{
	$y^{\prime}=-4\cdot 3\cdot x^2+\dfrac{1}{2}\cdot 2\cdot x-2+0=-12 x^2+x-2$
		\begin{itemchoice}
		\itemch Sai. Vì $a+b+c=-12$.
		\itemch Sai. Phương trình $ y'=0$ vô  nghiệm.
		\itemch Đúng. $y'(0) = -2$, nên đồ thị hàm số $ y'$ cắt trục tung tại điểm $\left(0;-2\right)$.
		\itemch Sai. Vì ta có $y'=3 \Leftrightarrow -12x^2+x-5=0$ vô nghiệm.
	\end{itemchoice}
	.}
\end{ex}

\begin{ex}%Câu 36%[1D7H2-3]
Cho hàm số $y=\dfrac{x-3}{2 x+1}$. Khi đó
\choiceTF
{\True $ y'(0)=7$}
{Đồ thị của hàm số $ y'$ đi qua điểm $ A\left(1;\dfrac{7}{3}\right)$}
{$ y'(1)<y'(2)$}
{\True Điểm $ M$ thuộc đồ thị $(C)$của hàm số $y=\dfrac{x-3}{2 x+1}$ có hoành độ $x_0=0$. Khi đó, phương trình tiếp tuyến của $(C)$ tại $M$ song song với đường thẳng $ y=7x+2024$}
\loigiai{
$y^{\prime}=\dfrac{(x-3)^{\prime}(2 x+1)-(2 x+1)^{\prime}(x-3)}{(2 x+1)^2}=\dfrac{2 x+1-2(x-3)}{(2 x+1)^2}=\dfrac{7}{(2 x+1)^2}$.
	\begin{itemchoice}
	\itemch Đúng. Vì $y'(0) =7$.
	\itemch Sai. Vì $y'(1)= \dfrac{7}{9} \ne  \dfrac{7}{3}$.
	\itemch Sai. $y'(1) = \dfrac{7}{9} > y'(2)=\dfrac{7}{25}$.
	\itemch Đúng. Vì ta có $y'(0)=7$.
\end{itemchoice}
}
\end{ex}

\begin{ex}%Câu 37%[1D7H2-1]
Tính được đạo hàm của các hàm số sau. Khi đó
\choiceTF
{$y=2\sin x-3\cos x$ có $y^{\prime}=2\cos x-3\sin x$}
{\True $y=3\cot x-\tan x$có $y^{\prime}=-\dfrac{2}{\sin^2x}-\dfrac{1}{\cos^2x}$}
{$y=x\cos x$ có $y^{\prime}=\cos x+x\sin x$}
{\True $y=2 x\sin ^2 x$ có $y^{\prime}=2\sin^2x+2x\sin 2x$}
\loigiai{
	\begin{itemchoice}		
	\itemch  $y^{\prime}=2(\sin x)^{\prime}-3(\cos x)^{\prime}=2\cos x+3\sin x$.
	\itemch $y^{\prime}=2(\cot x)^{\prime}-(\tan x)^{\prime}=-\dfrac{2}{\sin ^2 x}-\dfrac{1}{\cos ^2 x}$.\\
	\itemch $y^{\prime}=x^{\prime}\cdot\cos x+(\cos x)^{\prime}\cdot x=\cos x-x\sin x$.\\
	\itemch Ta có\\
	$\begin{aligned}
y^{\prime}&=(2 x)^{\prime}\cdot\sin ^2 x+\left(\sin ^2 x\right)^{\prime}\cdot 2 x=2\sin ^2 x+2\sin x(\sin x)^{\prime}\cdot 2 x\\
&=2\sin ^2 x+4 x\sin x\cos x=2\sin ^2 x+2 x\sin 2 x.
\end{aligned}$
\end{itemchoice}	
}
\end{ex}

\begin{ex}%Câu 38%[1D7V3-3]
Một vật chuyển động trên đường thẳng được xác định bởi công thức $s(t)=t^3-3 t^2+7 t-2$, trong đó $t>0$ và tính bằng giây và $s$ là quãng đường chuyển động được của vật trong $t$ giây tính bằng mét. Khi đó
\choiceTF
{\True Tốc độ của vật tại thời điểm $t=2$ là $v(2)=s^{\prime}(2)=7$ m/s}
{\True Gia tốc của vật tại thời điểm $t=2$ là $a(2)=v^{\prime}(2)=6$ m/s$^2$}
{Gia tốc của vật tại thời điểm mà vận tốc của chuyển động bằng $16$ m/s$^2$ là $ 10$ m/s$^2$}
{\True Thời điểm $t=1$ giây tại đó vận tốc của chuyển động đạt giá trị nhỏ nhất}
\loigiai{
Ta có $s^{\prime}(t)=3 t^2-6 t+7$ và $s^{\prime\prime}(t)=6 t-6$.
\begin{itemchoice}		
	\itemch  Vận tốc của vật tại thời điểm $t=2$ là $v(2)=s^{\prime}(2)=3\cdot 2^2-6\cdot 2+7=7$ (m/s).
	\itemch  Gia tốc của vật tại thời điểm $t=2$ là $a(2)=v^{\prime}(2)=s^{\prime\prime}(2)=6\cdot 2-6=6$ (m/s$^2$).
	\itemch Vận tốc của chuyển động bằng $16$ m/s$^2$ tại thời điểm $t$ nghĩa là\\	
$ v(t)=s^{\prime}(t)=16 \Leftrightarrow 3t^2-6t+7=16\Leftrightarrow \hoac{& t=3\text{ (thoả mãn)} \\ 
& t=-1\text{(loại).}}$\\
Gia tốc của vật tại thời điểm $t=3$ là $a(3)=v^{\prime}(3)=s^{\prime\prime}(3)=6\cdot 3-6=12$ (m/s$^2$).
	\itemch Vận tốc của chuyển động có phương trình $v(t)=3 t^2-6 t+7$ là một parabol, có đỉnh $S\left(-\dfrac{b}{2 a};-\dfrac{\Delta}{4 a}\right)\Rightarrow S(1 ; 4)$ và hệ số $a=3>0$ nên hàm số có giá trị nhỏ nhất bằng $4$ tại $t=1$. Vậy tại thời điểm $t=1$ thì vận tốc của chuyển động đạt giá trị nhỏ nhất bằng $4$ m/s.
\end{itemchoice}
}
\end{ex}





	\Closesolutionfile{ans}
	\indapan{3}{ans/ans-K11_C7_B2}
	
	\Opensolutionfile{ans}[ans/ans-K11_C7_B2]
	\caukq
	
	\begin{ex}%Câu 39
		Một chất điểm chuyển động theo phương trình $s(t)=3\sin 2 t+2\cos 2 t$, trong đó $t$ là thời gian tính bằng giây và $s$ là quãng đường chuyển động được của chất điểm trong $t$ giây tính bằng mét. Tính gia tốc của chất điểm đó khi $t=\dfrac{\pi}{4}$.
	\shortans{$-12$}
	\loigiai{
		Ta có $s^{\prime}(t)=3(\sin 2 t)^{\prime}+2(\cos 2 t)^{\prime}=6\cos 2 t-4\sin 2 t$.\\
		Và $s^{\prime\prime}(t)=6(\cos 2 t)^{\prime}-4(\sin 2 t)^{\prime}=-12\sin 2 t-8\cos 2 t$.\\
		Gia tốc của chất điểm tại thời điểm $t=\dfrac{\pi}{4}$ là \\
		$a\left(\dfrac{\pi}{4}\right)=s^{\prime\prime}\left(\dfrac{\pi}{4}\right)=-12\left[\sin\left(2\cdot\dfrac{\pi}{4}\right)\right]-8\left[\cos\left(2\cdot\dfrac{\pi}{4}\right)\right]=-12$.}
\end{ex}

\begin{ex}%Câu 40
Nhiệt độ cơ thể của một người trong thời gian bị bệnh được cho bởi công thức $T(t)=-0,1 t^2+1,2 t+98,6$, trong đó $T$ là nhiệt độ (tính theo đơn vị đo Fahrenheit) tại thời điểm $t$ (tính theo ngày). Tìm tốc độ thay đổi nhiệt độ ở thời điểm $t=2$.
\shortans{$ 0{,}8$}
\loigiai{
Đạo hàm của hàm số $T$ biểu thị tốc độ thay đổi của nhiệt độ.\\
Ta có $T^{\prime}(t)=-0{,}2 t+1{,}2$.\\
Vậy tốc độ thay đổi nhiệt độ tại thời điểm $t=2$ là
$T^{\prime}(2)=-0{,}2.2+1,2=0{,}8$ ($^\circ F$/ngày).
}
\end{ex}

\begin{ex}%Câu 41
Cho hàm số $y=-x^3+3 x^2+9 x-1$ có đồ thị là $(C)$. Tìm hệ số góc lớn nhất của tiếp tuyến tại một điểm $M$ trên đồ thị $(C)$.\\
\shortans{ $12$}
\loigiai{
Gọi điểm $M\left(x_0 ; y_0\right)\in(C)$ là toạ độ tiếp điểm và $y^{\prime}=-3 x^2+6 x+9$.\\
Hệ số góc của tiếp tuyến tại điểm $M$ là $y^{\prime}\left(x_o\right)=-3 x_o^2+6 x_o+9$.\\
Ta thấy, hệ số góc của tiếp tuyến tại điểm $M$ là một hàm số có đồ thị là một parabol, có đỉnh $S\left(-\dfrac{b}{2 a};-\dfrac{\Delta}{4 a}\right)\Rightarrow S(1 ; 12)$ và hệ số $a=-3<0$ nên hàm số có giá trị lớn nhất bằng $12$ tại $x_o=1$.\\
Vậy hệ số góc lớn nhất của tiếp tuyến là $y^{\prime}(1)=12$.
}
\end{ex}

\begin{ex}%Câu 42
Cho hàm số $y=f(x)=\sqrt{\tan x+\cot x}$. Tính $f^{\prime}\left(\dfrac{\pi}{4}\right)$.\\
\shortans{$0$}
\loigiai{
$f^{\prime}(x)=\dfrac{1+\tan ^2 x-\left(1+\cot ^2 x\right)}{2\sqrt{\tan x+\cot x}}=\dfrac{\tan ^2 x-\cot ^2 x}{2\sqrt{\tan x+\cot x}}\Rightarrow f^{\prime}(0)=\dfrac{\tan ^2\dfrac{\pi}{4}-\cot ^2\dfrac{\pi}{4}}{2\sqrt{\tan\dfrac{\pi}{4}+\cot\dfrac{\pi}{4}}}=0$.}
\end{ex}

\begin{ex}%Câu 43
Một chất điểm chuyển động theo phương trình $s(t)=10+t+9 t^2-t^3$ trong đó $s$ tính bằng mét, $t$ tính bằng giây. Tính thời gian để vận tốc của chất điểm đạt giá trị lớn nhất (tính từ thời điểm ban đầu)?\\
\shortans{$1{,}5$}
\loigiai{
Ta có $v(t)=s^{\prime}(t)=-3 t^2+9 t+1$ có đồ thị là Parabol, do đó $v(t)_{\max}\Leftrightarrow t=\dfrac{-9}{-6}=\dfrac{3}{2}=1{,}5$.}
\end{ex}

\begin{ex}%Câu 44
Một vật chuyển động với phương trình $S(t)=4 t^2+t^3$, trong đó $t>0, t$ tính bằng giây, $S(t)$ tính bằng $m$. Tìm gia tốc của vật tại thời điểm vận tốc của vật bằng 11 .\\
\shortans{$ 14$}
\loigiai{
Ta có $v(t)=S^{\prime}(t)=3 t^2+8 t$ và $a(t)=v^{\prime}(t)=6 t+8$.\\
Theo đề bài, ta có $v=11\Rightarrow 3 t^2+8 t=11\Rightarrow t=1$ $(t>0)$.\\
Vậy gia tốc của vật tại thời điểm vận tốc của vật bằng 11 là $a(1)=14$ m/s$^2$.}
\end{ex}
	








	
	\Closesolutionfile{ans}
	\indapan{6}{ans/ans-K11_C7_B2}